Como trabajo futuro se recomienda adaptar la amplitud del eco a cada ventana, usando más eco donde la música lo enmascara y menos donde se oiría, para mejorar a la vez la calidad y la recuperación; evaluar el sistema con archivos comprimidos, con ruido o remuestreados; reemplazar la repetición por un código corrector de errores con menos redundancia, que también funcione en audios cortos; y medir la calidad con una métrica diseñada para audio, como la evaluación perceptual de la calidad de audio (PEAQ)~\cite{itu_bs1387}, en lugar de PESQ, pensada para voz.
